% The phenomenon: LLMs finetuned on time series match or beat domain-specific models.
% The question: Why should representations learned from text help with temporal data?
% We provide a concrete, mechanistic account.
%
% Our approach: investigate from first principles.
%   1. Establish that transfer actually happens and quantify it. -> raw benchmark perf of pretrain Vs random
%   2. Characterize what the pretrained model already has (what does it contain that is useful for time series?) Analysis on QWEN BASE as it pertains to transfer 
%   3. Characterize what changes during finetuning (what changes in the representation space?) Analsysi on QWEN finetuned Compare with the control (learning from scratch) comparison of QWEN finetuend w random
%   4. Conclusion on how transfer works 
%
% Preview the main finding: pretraining creates a geometric prior —
% a low-dimensional, smooth representation space with directions
% that naturally encode temporal dynamics. Transfer = selecting those directions.

\section{Introduction}
\label{sec:intro}

% The key contributions we make in this paper are as follows:
% \begin{itemize}[leftmargin=*]
%   \item TODO
% \end{itemize}